\section{The actual arithmetic observable under spatial diffusion}
\label{sec:eo}

Retain the original real spatial coordinates $q=(x,y,w)$ and the
polynomial $F$ in \eqref{hc:F}. The complex formulas below are their
holomorphic extensions; a sum of derivative squares is a Euclidean
squared norm only on the stated real locus. Define the exact observable
\[
\sigma(q)=F_1(q)/2,\qquad \Psi(t,q)=Z_t(\sigma(q)).
\]
The Newman time $t$ and the spatial flow time $\tau$ remain distinct
variables. The already proved entire function $Z$ makes $\Psi$ entire
in $(t,q)\in\C\times\C^3$.

\subsection{Original vector fields and the rank-one heat operator}

Here is a direct determinant proof needed to define the vector fields.
Use the already proved inverse-chart map
\[
H(\alpha,r,c)=(\alpha^{-1},r-\alpha,
                  5\alpha^2-3r\alpha-c\alpha^3),\quad \alpha\ne0.
\]
Its inverse on $x\ne0$ is $(1/x,y+1/x,F_3(q))$.
The two derivative matrices, in the displayed coordinate order, are
\[
DH=\begin{pmatrix}
-\alpha^{-2}&0&0\\
-1&1&0\\
-3c\alpha^2+10\alpha-3r&-3\alpha&-\alpha^3
\end{pmatrix},\quad
D(F\circ H)=\begin{pmatrix}
r&\alpha-3cr^2+2r&-r^3\\
2&4-6cr&-3r^2\\
0&0&1
\end{pmatrix}.
\]
The first determinant is $\alpha$. Expansion along the last row of
the second gives $r(4-6cr)-2(\alpha-3cr^2+2r)=-2\alpha$.
The chain rule therefore proves $\det DF=-2$ on $x\ne0$.
Both sides are polynomials in $(x,y,w)$, so their equality extends
to $x=0$ by continuity. Put $J=DF$, $W_i=J^{-1}e_i$, and
\[
U=2W_1+6F_3W_2.
\]
These are globally defined polynomial vector fields, since
$J^{-1}=-(1/2)\operatorname{adj}J$. With $\epsilon_{123}=1$ and
cyclic indices, the adjugate formula is
$W_i=-\tfrac12\nabla F_{i+1}\times\nabla F_{i+2}$.
The divergence of each cross product is zero: its differentiated
components contract a symmetric Hessian with an antisymmetric
$\epsilon$ tensor. Also $W_iF_j=\delta_{ij}$, by $JW_i=e_i$.
The product rule now gives
\[
\operatorname{div}U=2\operatorname{div}W_1+
  6F_3\operatorname{div}W_2+6W_2F_3=0,\qquad U\sigma=1.
\]
For every twice differentiable scalar $h(s)$, direct differentiation
proves $U(h\circ\sigma)=h'(\sigma)$ and
$U^2(h\circ\sigma)=h''(\sigma)$. Consequently the actual observable
satisfies the exact operator identity
\begin{equation}
\partial_t\Psi+\tfrac14 U^2\Psi=0.
\label{eq:eo-rankone}
\end{equation}
Here the full spatial operator is
\[
U^2h=U^iU^j\partial_i\partial_jh+
             U^i(\partial_iU^j)\partial_jh
     =\partial_i(U^iU^j\partial_jh),
\]
where repeated indices run over the three original coordinates.
The last equality uses the proved divergence identity. Its real
principal tensor is $UU^{\mathsf T}/4$, of rank one because
$U\sigma=1$ implies $U\ne0$. On $C_c^\infty(\R^3)$, one integration
by parts gives $-\int\overline h U^2h=\int|Uh|^2$; polynomial
coefficients are bounded on the support and there are no boundary
terms. The explicit reversal $\Phi(\tau,q)=\Psi(-\tau,q)$ obeys
$\partial_\tau\Phi=U^2\Phi/4$. This proves an exact forward
spatial heat realization on the specified observable, with every
first-order coefficient retained.

\subsection{Every Euclidean diffusion coefficient}

Let $Q=1+xy$, an abbreviation for the original polynomial factor.
Differentiation before restriction to a fibre gives
\begin{align*}
g_1:=\partial_x\sigma&=\tfrac32 yQ^2w+
                           \tfrac12y^3(7+6xy),\\
g_2:=\partial_y\sigma&=\tfrac32 xQ^2w+
                           \tfrac12(8y+21xy^2+12x^2y^3),\\
g_3:=\partial_w\sigma&=\tfrac12Q^3.
\end{align*}
Define $N=g_1^2+g_2^2+g_3^2$ and
\[
L=\Delta_q\sigma=3wQ(x^2+y^2)+3y^4+18x^2y^2+21xy+4.
\]
For example the second derivatives in the Laplacian are
$\sigma_{xx}=3wy^2Q+3y^4$,
$\sigma_{yy}=3wx^2Q+18x^2y^2+21xy+4$, and
$\sigma_{ww}=0$, proving every term in $L$.
The chain rule in each coordinate, followed by summation, proves
\begin{equation}
\Delta_q\Psi=N(q)Z_{ss}(t,\sigma(q))+
                    L(q)Z_s(t,\sigma(q)).
\label{eq:eo-euclidean}
\end{equation}
No term is suppressed by the one-variable presentation of $Z$.
For real $q$, $N=|\nabla\sigma|^2>0$: an invertible $J$ cannot
have its first row zero.

On the original finite inverse lift
$\Gamma(r)=(-1/(2r),3r,26r^2)$, $r\ne0$, direct substitution gives
\[
\nabla\sigma(\Gamma(r))=
  (9r^3/4,3r/8,-1/16)^{\mathsf T},\qquad
D^2\sigma(\Gamma(r))=
\begin{pmatrix}
-108r^4&3r^2/4&9r/8\\
3r^2/4&13/4&-3/(16r)\\
9r/8&-3/(16r)&0
\end{pmatrix}.
\]
Thus, retaining $\sigma(\Gamma(r))=-r^2/2$,
\begin{align}
N(\Gamma(r))&=\frac{1296r^6+36r^2+1}{256},\notag\\
L(\Gamma(r))&=\frac{13}{4}-108r^4,\notag\\
(\Delta_q\Psi)(t,\Gamma(r))
 &=\frac{1296r^6+36r^2+1}{256}Z_{ss}(t,-r^2/2)
  +\left(\frac{13}{4}-108r^4\right)Z_s(t,-r^2/2).
\label{eq:eo-on-lift}
\end{align}
These are identities of rational holomorphic functions for complex
$r\ne0$. Their Euclidean-norm interpretation uses real $r\ne0$.

For a real differentiable time map $t=\theta(\tau)$ and the original
viscosity $\nu>0$, the complete advection--diffusion residual is
\begin{align}
&(\partial_\tau+U\cdot\nabla_q-\nu\Delta_q)
             \Psi(\theta(\tau),q)\notag\\
&\qquad=\left(-\frac{\theta'(\tau)}4-\nu N(q)\right)
                  Z_{ss}(\theta(\tau),\sigma(q))
       +(1-\nu L(q))Z_s(\theta(\tau),\sigma(q)).
\label{eq:eo-full-residual}
\end{align}
The three preceding exact identities prove this formula term by term.
It is the direct map between the actual arithmetic heat solution and
the scalar spatial advection--diffusion operator of this polynomial
flow. It does not assert that its displayed residual vanishes.

\subsection{The particular arithmetic covector along the escaping orbit}

Local existence and uniqueness of the smooth polynomial ODE
$\partial_\tau X_\tau=U(X_\tau)$ follow on a closed finite ball
from the contraction map $Y\mapsto q+\int_0^\tau U(Y(v))\,dv$
for sufficiently short time; bounded $DU$ supplies its Lipschitz
constant. Continuing these solutions defines the maximal local flow
on its actual open domains. Smooth dependence follows by differentiating
the integral equation; uniqueness supplies the inverse local flow.
The identities $UF_1=2$, $UF_2=6F_3$, and $UF_3=0$ integrate to
\[
F(X_\tau(q))=(F_1(q)+2\tau,F_2(q)+6\tau F_3(q),F_3(q)).
\]
In particular $\sigma(X_\tau(q))=\sigma(q)+\tau$. Write
$A_\tau=D_qX_\tau$ and $K_\tau=A_\tau^{-1}A_\tau^{-\mathsf T}$.
Differentiating this last scalar identity gives the full covector map
\[
A_\tau^{\mathsf T}\nabla\sigma(X_\tau(q))=\nabla\sigma(q),\qquad
\nu\nabla\sigma(q)^{\mathsf T}K_\tau\nabla\sigma(q)
       =\nu N(X_\tau(q)).
\]
Every matrix is evaluated at the same initial point $q$ and time
$\tau$; no arbitrary covector has replaced this actual one.

At $q_0=(1,-3/2,13/2)=\Gamma(-1/2)$ put
$z=\sqrt{1-8\tau}>0$, with inverse $\tau=(1-z^2)/8$.
For $-\infty<\tau<1/8$ the explicit trajectory is
$X_\tau(q_0)=\Gamma(-z/2)=(z^{-1},-3z/2,13z^2/2)$.
To verify the ODE, its $F$ image is $(-1/4+2\tau,0,0)$;
differentiation gives $J\gamma'=(2,0,0)=JU(\gamma)$,
so invertibility of $J$ gives $\gamma'=U(\gamma)$. Its value at
$\tau=0$ is $q_0$, and its escape as $\tau\uparrow1/8$ precludes
a finite continuous extension of the ODE solution. For the specific
initial covector
\[
\ell_\sigma=\nabla\sigma(q_0)=(-9/32,-3/16,-1/16)^{\mathsf T}
\]
we have the entire transported vector and cost
\begin{align*}
A_\tau^{-\mathsf T}\ell_\sigma
 &=(-9z^3/32,-3z/16,-1/16)^{\mathsf T},\\
\nu\ell_\sigma^{\mathsf T}K_\tau\ell_\sigma
 &=\nu\frac{81z^6+36z^2+4}{1024}
   \longrightarrow\frac{\nu}{256}\quad(z\downarrow0).
\end{align*}
This computes the precise diffusive cost of the covector selected
by the arithmetic observable, even at the spatial escape endpoint.
The zero condition itself remains $Z_t(\sigma(q))=0$ with its full
analytic divisor as proved above; the transported cost is not an
assertion of such a zero or of a Euclidean Navier--Stokes counterexample.
